\begin{table}[t]
\centering
\footnotesize
\setlength{\tabcolsep}{6pt}
\begin{threeparttable}
\caption{%
\textbf{Scaling the visual and text encoders.}
Larger ViT backbones improve downstream performance, with the strongest gains on tasks requiring fine-grained text understanding (PODS). Utilizing RoBERTa-Base instead of Large causes modest drops in feature preservation (FG-CLS, ADE20k) but does not negatively affect \model{}'s text understanding. 
}
\label{tab:scaling}
\begin{tabular}{@{} l l
                S[table-format=2.1]
                S[table-format=2.1]
                S[table-format=2.1]
                S[table-format=2.1]@{}}
\toprule
{Component} & {Scale}
  & {FG-CLS $\uparrow$}
  & {ADE20k $\uparrow$}
  & {CORE $\uparrow$}
  & {PODS $\uparrow$} \\
\midrule
DINOv2 & ViT-S & {74.4} & {50.4} & {91.8} & {45.5} \\
\rowcolor{\methodcolor}
DINOv2 & ViT-B & \textbf{76.2} & {55.1} & {94.9} & {50.6} \\
DINOv2 & ViT-L & {74.8}  & \textbf{56.2} & \textbf{96.2} & {53.0} \\
\addlinespace
% 4 & SigLIP & ViT-B & 64.3 & 90.9 & {--} & 23.0 \\
% 5 & SigLIP & ViT-L & {XX.X} & {XX.X} & {XX.X} & {XX.X} \\
% \addlinespace
RoBERTa & Base & {{72.7}}  & {54.8} & {94.8} &  {\textbf{53.9}} \\
\rowcolor{\methodcolor}
% 5 & RoBERTa & \textbf{Large (ours)} & \textbf{77.0} & \textbf{93.4} & \textbf{55.3} & \textbf{55.3} \\
\bottomrule
\end{tabular}
\begin{tablenotes}[flushleft]\footnotesize
\scriptsize
\item ViT scaling w/ RoBERTa-Large; RoBERTa scaling w/ DINOv2 ViT-B.
% All models trained for 500k iterations on the full data mixture.
% PODS uses descriptive text prompts.
\end{tablenotes}
\end{threeparttable}
\end{table}
